% !TEX program = pdflatex
% Boson14 Planted-Clique Benchmark Suite — Presentation Manual
% Compile: pdflatex manual.tex (run twice for correct references)

\documentclass[aspectratio=169]{beamer}
\usetheme{metropolis}

\usepackage{amsmath,amssymb}
\usepackage{booktabs}
\usepackage{listings}
\usepackage{xcolor}
\usepackage{graphicx}
\usepackage{hyperref}

% Listings style for code blocks
\lstset{
    basicstyle=\ttfamily\footnotesize,
    backgroundcolor=\color{black!5},
    frame=single,
    framerule=0pt,
    breaklines=true,
    breakatwhitespace=true,
    columns=fullflexible,
    keepspaces=true,
    showstringspaces=false,
    keywordstyle=\color{blue!70!black},
    commentstyle=\color{green!50!black},
    stringstyle=\color{orange!80!black},
    language=Python,
}

\lstdefinestyle{shell}{
    language=bash,
    keywordstyle=\color{blue!70!black},
    morekeywords={python,uv,pip,cd,export},
}

% Metadata
\title{Boson14 Planted-Clique Benchmark Suite}
\subtitle{Scaled Motzkin-Straus Calibration for Dirac-3}
\author{QCi Hardware Team}
\date{\today}
\institute{Quantum Computing Inc.}

\begin{document}

% ── Title ─────────────────────────────────────────────────────────────────

\maketitle

% ── Why Planted Cliques? ─────────────────────────────────────────────────

\begin{frame}{Why Planted Cliques?}
    \begin{itemize}
        \item \textbf{Known $\omega$} — ground truth is planted, not discovered
        \item \textbf{Controllable difficulty} — adjust $n$, $k$, edge probability $p$
        \item \textbf{Integer coupling} — $J = -A$ has entries in $\{0, -1\}$, no half-integer scaling
        \item \textbf{Degenerate landscapes} — plant multiple same-size cliques to test multi-modal behaviour
        \item \textbf{Blind testing} — scramble vertex labels; solver sees no hint of planted structure
    \end{itemize}

    \vspace{1em}
    \alert{Key insight:} If hardware recovers the planted $\omega$ across different $R$ values and permutations, calibration is verified end-to-end.
\end{frame}

% ── Motzkin-Straus Theorem ───────────────────────────────────────────────

\begin{frame}{Motzkin-Straus Theorem}
    The clique number can be found via a continuous optimisation:

    \begin{theorem}[Motzkin \& Straus, 1965]
        \[
            \max_{x \in \Delta_n} \frac{1}{2} x^\top A x = \frac{1}{2}\left(1 - \frac{1}{\omega}\right)
        \]
        where $A$ is the adjacency matrix and $\Delta_n = \{x \geq 0 : \sum x_i = 1\}$ is the standard simplex.
    \end{theorem}

    \vspace{0.5em}
    \textbf{Inverting:} Given optimal objective $f^*$:
    \[
        \omega = \operatorname{round}\!\left(\frac{1}{1 - 2f^*}\right)
    \]
\end{frame}

% ── Scaled Formulation ───────────────────────────────────────────────────

\begin{frame}{Scaled Simplex Formulation}
    Scale to integer-friendly constraints: $y = R\,x$, so $\sum y_i = R$.

    \begin{align*}
        g(y) &= \tfrac{1}{2}\, y^\top A\, y = R^2 \cdot f(x) \\[6pt]
        g^*(\omega) &= \frac{R^2}{2}\Bigl(1 - \frac{1}{\omega}\Bigr) \\[6pt]
        \omega &= \operatorname{round}\!\Bigl(\frac{R^2}{R^2 - 2\,g^*}\Bigr)
    \end{align*}

    \vspace{0.5em}
    \begin{itemize}
        \item Optimal allocation: $y_i = R/\omega$ for clique vertices, $0$ otherwise
        \item Coupling matrix: $J = -A$ (entries in $\{0, -1\}$, all integers)
    \end{itemize}
\end{frame}

% ── Hardware Energy Mapping ──────────────────────────────────────────────

\begin{frame}{Hardware Energy Mapping}
    Dirac-3 \textbf{minimises} $E(y) = y^\top J\, y$ on the scaled simplex $\sum y_i = R$.

    \begin{align*}
        E(y) &= y^\top J\, y = -y^\top A\, y = -2\,g(y) \\[6pt]
        E^* &= -R^2\Bigl(1 - \frac{1}{\omega}\Bigr) \\[6pt]
        \omega &= \operatorname{round}\!\Bigl(\frac{R^2}{R^2 + E^*}\Bigr)
    \end{align*}

    \vspace{0.5em}
    \begin{itemize}
        \item Both $g \to \omega$ and $E \to \omega$ inversions should agree on every sample
        \item Consistency check: count agreement rate across all Dirac-3 samples
    \end{itemize}
\end{frame}

% ── Toolkit File Structure ───────────────────────────────────────────────

\begin{frame}[fragile]{Toolkit File Structure}
\begin{lstlisting}[style=shell,numbers=none]
boson14/
  generate.py               # Instance generation (planted, degenerate, scramble)
  test_dirac.py             # Dirac-3 submission + verification
  scaled-motzkin-straus.md  # Mathematical derivation
  USAGE.md                  # Full CLI reference
  manual.pdf                # This presentation
  inputs/                   # External CSV files (e.g. DIMACS graphs)
    C125.9.csv              # 125-vertex dense graph (omega=34)
  output/                   # Results directory
    *_meta.json             # Instance metadata
    *_boson14.csv           # StQP adjacency CSV (n x n+1)
    *_dirac_meta.json       # Solver results
    *_solutions.npz         # Theoretical solutions
    *_dirac_solutions.npz   # Hardware samples
    *.png                   # Plots
\end{lstlisting}

    \vspace{0.5em}
    \textbf{Library:} Core functions live in \texttt{dirac\_bench.boson14} (the \texttt{dirac-test-suite} package).
\end{frame}

% ── generate.py Demo ─────────────────────────────────────────────────────

\begin{frame}[fragile]{generate.py --- Generating Instances}
\begin{lstlisting}[style=shell,numbers=none]
uv run python boson14/generate.py --n 30 --k 12 --verify
\end{lstlisting}

\begin{lstlisting}[numbers=none,basicstyle=\ttfamily\scriptsize]
Generating: n=30, k=12, p=0.5, seed=42, R=100
  Planted nodes: [2, 3, 6, 10, 14, 16, 19, 21, 23, 28, 29, 30]

Theoretical values (R=100, w=12):
  g* = 4583.33
  E* = -9166.67
  y_i = R/w = 100/12 = 8.3333 for planted vertices

Integer QP: J has entries in {-1, 0}

Brute-force verification:
  Clique number (brute-force): w = 12
  CONFIRMED: w = k = 12
\end{lstlisting}

    Variants:
\begin{lstlisting}[style=shell,numbers=none]
# Dense edges, custom seed
generate.py --n 30 --k 12 --density dense --seed 99
# With distribution plots
generate.py --n 40 --k 15 --plot-distribution --verify
\end{lstlisting}
\end{frame}

% ── Degenerate Cliques ───────────────────────────────────────────────────

\begin{frame}[fragile]{Degenerate Cliques}
    Plant multiple same-size cliques to create a \textbf{degenerate} optimisation landscape:

    \begin{itemize}
        \item All planted cliques have size $k$ --- same $\omega$, multiple optima
        \item Cross-edges between planted sets are \textbf{severed} to guarantee independence
        \item Tests whether hardware explores multiple basins or collapses to one
    \end{itemize}

    \vspace{0.5em}
\begin{lstlisting}[style=shell,numbers=none]
uv run python boson14/generate.py --n 50 --k 10 --num-cliques 3 --verify
\end{lstlisting}

\begin{lstlisting}[numbers=none,basicstyle=\ttfamily\scriptsize]
  Degenerate: 3 planted 10-cliques
  Set 1: [6, 18, 21, 26, 28, 40, 46, 47, 48, 50]
  Set 2: [5, 8, 19, 24, 25, 30, 32, 38, 39, 41]
  Set 3: [10, 16, 17, 22, 27, 29, 33, 35, 42, 43]
  omega_verified: 10, num_max_cliques: 3
\end{lstlisting}
\end{frame}

% ── Index Scrambling ─────────────────────────────────────────────────────

\begin{frame}[fragile]{Index Scrambling}
    Apply a random vertex permutation for blind testing:

    \begin{itemize}
        \item \textbf{Forward permutation:} original label $\to$ scrambled label
        \item \textbf{Inverse permutation:} scrambled label $\to$ original label (stored in metadata)
        \item \texttt{unscramble\_solution()} maps Dirac-3 y-vectors back to original ordering
        \item Objectives are invariant: $g(y_{\text{scrambled}}) = g(y_{\text{original}})$
    \end{itemize}

    \vspace{0.5em}
\begin{lstlisting}[style=shell,numbers=none]
# Generate scrambled
uv run python boson14/generate.py --n 50 --k 15 --num-cliques 3 --scramble
# Test scramble roundtrip on Dirac-3
uv run python boson14/test_dirac.py --n 14 --k 7 --scramble --verify --plot
\end{lstlisting}

\begin{lstlisting}[numbers=none,basicstyle=\ttfamily\scriptsize]
--- Scramble Roundtrip ---
  Sample 0: g(scrambled)=4285.65, sum(y_orig)=100.00
  Scramble roundtrip: objectives preserved through permutation
\end{lstlisting}
\end{frame}

% ── test_dirac.py Demo ───────────────────────────────────────────────────

\begin{frame}[fragile]{test\_dirac.py --- Dirac-3 Submission}
    Three input modes: inline parameters, JSON from \texttt{generate.py}, or CSV file:
\begin{lstlisting}[style=shell,numbers=none]
# Inline planted-clique parameters
uv run python boson14/test_dirac.py --n 14 --k 7 --verify --plot
# From generate.py JSON
uv run python boson14/test_dirac.py --from-json .../boson14.json --verify
# From CSV (planted or external graph)
uv run python boson14/test_dirac.py --csv .../boson14.csv --k 7 --verify
\end{lstlisting}

\begin{lstlisting}[numbers=none,basicstyle=\ttfamily\scriptsize]
============================================================
Boson14 Dirac-3 Test: n=14, k=7, p=0.5, seed=42, R=100
============================================================
Planted nodes: [3, 5, 6, 7, 8, 12, 13]

Theoretical values (R=100, w=7):
  g* = 4285.71,  E* = -8571.43,  y_i = 14.2857

--- Dirac-3 Submission (R=100) ---
  Submitting 14-variable scaled QP to Dirac-3 cloud
  Best g = 4285.65  =>  omega = 7  (119.5s)

============================================================
SUMMARY: Best omega = 7, Known omega = 7  =>  MATCH
============================================================
\end{lstlisting}
\end{frame}

% ── CSV Input Path ──────────────────────────────────────────────────────

\begin{frame}[fragile]{CSV Input Path}
    Load any graph from a CSV adjacency file --- no planted-clique generation required.

    \vspace{0.3em}
    \textbf{CSV format:} $n$ rows $\times$ $(n{+}1)$ columns. Column~0 = linear coefficients $C_i$ (all zeros for Motzkin-Straus), columns 1\ldots$n$ = symmetric adjacency $A$.

    \vspace{0.5em}
    \textbf{With known $k$} (planted-clique CSV from \texttt{generate.py}):
\begin{lstlisting}[style=shell,numbers=none]
uv run python boson14/test_dirac.py \
  --csv boson14/output/n14_k7_p0.5_s42/n14_k7_p0.5_s42_boson14.csv \
  --k 7 --verify --plot
\end{lstlisting}

    \textbf{External graph, unknown $k$} (e.g.\ DIMACS):
\begin{lstlisting}[style=shell,numbers=none]
uv run python boson14/test_dirac.py \
  --csv boson14/inputs/C125.9.csv --verify
\end{lstlisting}

    \vspace{0.3em}
    \begin{itemize}
        \item CSV$\to$JSON round-trip is verified automatically (adjacency survives conversion)
        \item Omit \texttt{--k} for unknown instances --- brute-force still reports true $\omega$
        \item Output saved to \texttt{boson14/output/\{csv\_stem\}/}
    \end{itemize}
\end{frame}

% ── CSV Sample Output ───────────────────────────────────────────────────

\begin{frame}[fragile]{CSV Input --- Sample Output}
\begin{lstlisting}[numbers=none,basicstyle=\ttfamily\scriptsize]
Loaded CSV: .../n14_k7_p0.5_s42_boson14.csv (n=14)
  CSV->JSON round-trip: PASS

============================================================
Boson14 Dirac-3 Test: n=14, k=7, R=100  (from CSV)
============================================================

Theoretical values (R=100, w=7):
  g* = 4285.71,  E* = -8571.43,  y_i = 14.2857

Brute-force verification:
  Clique number (brute-force): w = 7
  CONFIRMED: w = k = 7

--- Dirac-3 Submission (R=100) ---
  Best g = 4285.65  =>  omega = 7  (161.4s)

Conversion consistency:
  omega_from_g == omega_from_E: 100/100 samples agree

============================================================
SUMMARY: Best omega = 7, Known omega = 7  =>  MATCH
============================================================
\end{lstlisting}
\end{frame}

% ── Multi-R Sweep ────────────────────────────────────────────────────────

\begin{frame}[fragile]{Multi-R Sweep}
    Verify that omega recovery is consistent across different $R$ values:

\begin{lstlisting}[style=shell,numbers=none]
uv run python boson14/test_dirac.py --n 14 --k 7 --multi-R --verify
\end{lstlisting}

\begin{lstlisting}[numbers=none,basicstyle=\ttfamily\scriptsize]
--- Multi-R Sweep ---

      R |         g* |         E* |     best_g | best_omega
  -----+-----------+-----------+-----------+-----------
      1 |     0.4286 |    -0.8571 |     0.4286 |          7
     10 |    42.8571 |   -85.7143 |    42.8571 |          7
    100 |  4285.7143 | -8571.4286 |  4285.6501 |          7

  PASS: All R values recover omega = 7
\end{lstlisting}

    \vspace{0.5em}
    \begin{itemize}
        \item Tests $R \in \{1, 10, 100\}$ on the same graph instance
        \item All three must agree on $\omega$ for the test to PASS
        \item Validates that energy scaling does not affect clique recovery
    \end{itemize}
\end{frame}

% ── Generating Plots ─────────────────────────────────────────────────────

\begin{frame}[fragile]{Generating Distribution Plots}
    Two plot types are available:

    \vspace{0.3em}
    \begin{description}
        \item[Clique distribution] (\texttt{clique\_dist\_*.png}) — two-panel figure:
            top panel shows maximal clique counts by size,
            bottom panel shows the scaled objective $g^*$ for each maximal clique with
            vertical $\omega$ lines.
        \item[Scaled histogram] (\texttt{scaled\_histogram\_*.png}) — histogram of
            Dirac-3 sample objectives with theoretical $g^*(\omega)$ reference lines.
    \end{description}

    \vspace{0.3em}
    \textbf{From generate.py} (local brute-force, no hardware needed):
\begin{lstlisting}[style=shell,numbers=none]
uv run python boson14/generate.py --n 14 --k 7 --plot-distribution --verify
\end{lstlisting}

    \textbf{From test\_dirac.py} (both plot types, after Dirac-3 submission):
\begin{lstlisting}[style=shell,numbers=none]
uv run python boson14/test_dirac.py --n 14 --k 7 --verify --plot
\end{lstlisting}

    Plots are saved to \texttt{--output-dir} (default: \texttt{boson14/output/}).
\end{frame}

% ── Example Plot ─────────────────────────────────────────────────────────

\begin{frame}{Example: Clique Distribution ($n{=}14$, $k{=}7$)}
    \centering
    \includegraphics[width=0.78\textwidth]{output/n14_k7_p0.5_s42/clique_dist_n14_k7_p0.5_s42.png}

    \vspace{0.3em}
    {\small\texttt{generate.py --n 14 --k 7 --plot-distribution --verify}}
\end{frame}

% ── Scaled Histogram ─────────────────────────────────────────────────────

\begin{frame}{Scaled Histogram --- What It Shows}
    The histogram plots the distribution of $g(y) = \tfrac{1}{2}y^\top A y$ across all Dirac-3 samples:

    \vspace{0.5em}
    \begin{itemize}
        \item \textbf{Vertical lines} mark theoretical $g^*(\omega)$ for integer $\omega$ values
        \item \textbf{Green line:} known $\omega$ (planted clique size)
        \item \textbf{Red line:} best objective achieved by hardware
        \item \textbf{Histogram shape:} indicates consistency of hardware output
    \end{itemize}

    \vspace{0.5em}
    \alert{Tight clustering near the planted $\omega$ line = well-calibrated hardware.}

    \vspace{0.3em}
    \begin{displaymath}
        g^*(\omega) = \frac{R^2}{2}\Bigl(1 - \frac{1}{\omega}\Bigr)
    \end{displaymath}
\end{frame}

% ── Interpreting Results ─────────────────────────────────────────────────

\begin{frame}{Interpreting Results}
    \textbf{What histograms tell you about hardware quality:}

    \vspace{0.5em}
    \begin{description}
        \item[Peak at planted $\omega$:] Hardware finds optimal solutions consistently.
        \item[Peak 1--2 below planted $\omega$:] Good performance, slight calibration drift.
        \item[Wide spread:] Hardware output is noisy; may need parameter tuning.
        \item[Bimodal distribution:] Hardware may be finding multiple local optima (expected for degenerate instances).
    \end{description}

    \vspace{0.5em}
    \textbf{Tuning knobs:}
    \begin{itemize}
        \item \texttt{relaxation\_schedule}: 1--4, controls annealing behaviour
        \item \texttt{num\_samples}: More samples $\to$ better statistics
        \item \texttt{R}: Larger $R$ $\to$ larger energy gap between clique sizes
    \end{itemize}
\end{frame}

% ── Output Formats ───────────────────────────────────────────────────────

\begin{frame}[fragile]{Output Formats}
    \begin{tabular}{@{}lll@{}}
        \toprule
        \textbf{Suffix} & \textbf{Source} & \textbf{Contents} \\
        \midrule
        \texttt{\_meta.json}           & generate.py   & Graph params, planted sets, $g^*$, $E^*$ \\
        \texttt{\_dirac\_meta.json}    & test\_dirac.py & Solver params, best $\omega$, solve time \\
        \texttt{\_solutions.npz}       & both          & Theoretical max-clique y-vectors \\
        \texttt{\_dirac\_solutions.npz}& test\_dirac.py & Dirac-3 sample y-vectors \\
        \texttt{clique\_dist\_*.png}   & both          & Clique size distribution bar chart \\
        \texttt{scaled\_histogram\_*.png} & test\_dirac.py & Objective histogram \\
        \bottomrule
    \end{tabular}

    \vspace{0.8em}
    Loading solutions in Python:
\begin{lstlisting}[basicstyle=\ttfamily\scriptsize]
import numpy as np
data = np.load("boson14/output/n14_k7_p0.5_s42_dirac_solutions.npz")
y_vectors = data["y_vectors"]  # shape (num_samples, n)
\end{lstlisting}
\end{frame}

% ── Worked Example ───────────────────────────────────────────────────────

\begin{frame}{Worked Example ($R=100$, $\omega=7$)}
    \centering
    \begin{tabular}{@{}lcr@{}}
        \toprule
        \textbf{Quantity} & \textbf{Formula} & \textbf{Value} \\
        \midrule
        $f^*$              & $\frac{1}{2}(1 - 1/7)$                    & $0.4286$ \\[4pt]
        $g^*$              & $\frac{100^2}{2}(1 - 1/7)$               & $4285.71$ \\[4pt]
        $E^*$              & $-100^2(1 - 1/7)$                        & $-8571.43$ \\[4pt]
        $y_i$ (clique)     & $100/7$                                   & $14.2857$ \\[4pt]
        $y_i$ (non-clique) & $0$                                       & $0$ \\
        \bottomrule
    \end{tabular}

    \vspace{1em}
    \raggedright
    \textbf{Verification:}
    \begin{itemize}
        \item $\omega = \operatorname{round}(10000 / (10000 - 8571.43)) = \operatorname{round}(7) = 7$ \checkmark
        \item $\omega = \operatorname{round}(10000 / (10000 + (-8571.43))) = \operatorname{round}(7) = 7$ \checkmark
        \item $\sum y_i = 7 \times 14.2857 = 100 = R$ \checkmark
    \end{itemize}
\end{frame}

\end{document}
